\chapter{Konzeption und Anforderungen}
\label{analyse-und-konzeption}

In diesem Kapitel werden die Anforderungen an die zu entwickelnde Lösung in Form von 
User-Stories nach Cohn \cite{cohn_user_2004} definiert. Aus jeder User-Story werden 
anschließend die Aufgaben zu ihrer Umsetzung sowie Akzeptanzkriterien (Definition of Done) 
abgeleitet, anhand derer in Abschnitt~\ref{requirements-check} überprüft wird, ob die 
jeweilige Anforderung erfüllt ist.

\noindent Die User-Stories bilden dabei das Bindeglied zwischen den Zielen aus 
Kapitel~\ref{motivation-und-zielsetzung} und dem Endprodukt dieser Arbeit: Während die 
übergeordneten Ziele und die daraus abgeleiteten Arbeitsziele beschreiben, was in dieser 
Arbeit umgesetzt wird, beschreiben die User-Stories aus Sicht der späteren Nutzer, welchen 
Nutzen diese Umsetzung stiften soll. Das Endprodukt besteht aus der 
Prozessverwaltungsplattform, welche die Orchestrierung übernimmt, sowie dem Compiler, der 
die Ausführung und Codegenerierung ermöglicht. Zu jeder User-Story wird daher angegeben, 
welche Arbeitsziele zu ihrer Erfüllung beitragen.

\label{requirements}

\begin{enumerate}[label=\Alph*)]
    \item Ich als \textbf{Softwareentwickler} möchte \textbf{Produktionsanlagen einfach an die Prozessschicht anbinden},
    sodass \textbf{neue Anlagen integriert werden können}.

    \item Ich als \textbf{Geschäftsleiter} möchte \textbf{meine Produktionsprozesse gemäß meiner Modellierung durchführen},
    sodass \textbf{die Produktionen selbst konsistent und messbar sind}.

    \item Ich als \textbf{Controller} möchte \textbf{laufende Prozesse überwachen}, sodass \textbf{die Effizienz von Geschäftsprozessen 
    gemessen und optimiert werden kann}.

    \item Ich als \textbf{Analyst} möchte \textbf{einen Verlauf vergangener Prozesse einsehen}, sodass \textbf{Schwachstellen und Engpässe
    ermittelt werden können}.

    \item Ich als \textbf{Kunde} möchte \textbf{über den aktuellen Stand meines Produktes informiert werden}, sodass \textbf{ich die nächsten
    Schritte mit dem Produkt planen kann}.

    \item Ich als \textbf{Administrator} möchte \textbf{ein wartbares IoT-System haben}, sodass \textbf{im Falle eines Ausfalls schnell reagiert
    werden kann, um den Produktionsfluss nicht zu stark zu beeinträchtigen}.
\end{enumerate}


\noindent \textbf{User-Story A: Softwareentwickler}

\textbf{\textit{Zugehörige Arbeitsziele}}: AZ1, AZ2

\textbf{\textit{Abzuleitende Tasks}}:

\begin{itemize}
    \item Entwicklung eines \ac{SDK} zur Kommunikation mit der Prozessschicht auf Ebene der Anlage.
    \item Entwicklung eines Werkzeugs zur Überführung des Prozessmodells in Firmwarecode.
\end{itemize}

\textbf{\textit{Definitions of Done}}:

\begin{itemize}
    \item Aus einem EPK-Prozessmodell lässt sich eine Schablone zur Entwicklung von Firmwarecode erstellen.
    \item Die Schablone lässt sich fehlerfrei kompilieren.
\end{itemize}


\noindent \textbf{User-Story B: Geschäftsleiter}

\textbf{\textit{Zugehörige Arbeitsziele}}: AZ3, AZ4, AZ5, AZ7

\textbf{\textit{Abzuleitende Tasks}}:

\begin{itemize}
    \item Entwicklung eines Compilers zur Übersetzung des EPK-Modells in eine ausführbare Form.
    \item Entwicklung einer Execution-Engine zur eigentlichen Ausführung.
    \item Einbau von Logging zur Zeitmessung bei der Ausführung.
\end{itemize}

\textbf{\textit{Definitions of Done}}:

\begin{itemize}
    \item Ein EPK-Modell lässt sich in das Ausführungsformat übersetzen.
    \item Das Ausführungsformat wird gemäß dem EPK-Modell ausgeführt.
    \item Am Ende einer Ausführung steht ein Ausführungsbericht zur Verfügung.
\end{itemize}


\noindent \textbf{User-Story C: Controller}

\textbf{\textit{Zugehörige Arbeitsziele}}: AZ6

\textbf{\textit{Abzuleitende Tasks}}:

\begin{itemize}
    \item Entwicklung von Live-Monitoring für einen laufenden Prozess.
\end{itemize}

\textbf{\textit{Definitions of Done}}:

\begin{itemize}
    \item Die Ausführung eines Prozessdurchlaufs lässt sich in der Prozesssicht visuell mitverfolgen.
\end{itemize}


\noindent \textbf{User-Story D: Analyst}

\textbf{\textit{Zugehörige Arbeitsziele}}: AZ7

\textbf{\textit{Abzuleitende Tasks}}:

\begin{itemize}
    \item Entwicklung von Event-Logging, sodass einzelne Ausführungsschritte aufgezeichnet werden.
\end{itemize}

\textbf{\textit{Definitions of Done}}:

\begin{itemize}
    \item Nach einer Ausführung ist es möglich, einen Ausführungsbericht zu generieren.
    \item Dieser Ausführungsbericht liegt in einem Format vor, mit dem sich Engpässe ermitteln lassen.
\end{itemize}


\noindent \textbf{User-Story E: Kunde}

\textbf{\textit{Zugehörige Arbeitsziele}}: keine -- diese User-Story geht über den Umfang 
der Arbeitsziele hinaus und wird in dieser Arbeit nur konzeptionell betrachtet.

\textbf{\textit{Abzuleitende Tasks}}:

\begin{itemize}
    \item Entwicklung von Live-Monitoring für einen laufenden Prozess.
\end{itemize}

\textbf{\textit{Definitions of Done}}:

\begin{itemize}
    \item Die Ausführung eines Prozessdurchlaufs lässt sich für den Kunden über eine eigene Sicht visuell mitverfolgen.
    \item Es können Benachrichtigungen an den Kunden verschickt werden, wenn ein bestimmter Schritt ausgeführt wurde.
\end{itemize}


\noindent \textbf{User-Story F: Administrator}

\textbf{\textit{Zugehörige Arbeitsziele}}: AZ1, AZ2, AZ7

\textbf{\textit{Abzuleitende Tasks}}:

\begin{itemize}
    \item Entwicklung einer modularen Hardwarearchitektur.
    \item Auswahl eines Kommunikationsprotokolls zwischen Prozessgateway und Produktionshardware.
    \item Definition eines Datenformats für die Kommunikation zwischen Prozessgateway und Produktionshardware.
\end{itemize}

\textbf{\textit{Definitions of Done}}:

\begin{itemize}
    \item Die Prozesshardware lässt sich einfach austauschen, ohne das gesamte System zu beeinflussen.
    \item Anhand der Ausführungsberichte lässt sich fehlerhafte Hardware schnell ausfindig machen.
\end{itemize}